\section{常系数线性差分方程}

\subsection{常系数线性差分方程的经典解法}\label{sub:difference_equation}

对于离散系统，一般可以用\textbf{差分方程}来描述激励与响应的关系。

\begin{definition}[差分]
    对于离散信号$x[n]$，定义其\textbf{前向差分}为
    $$\Delta x[n]=x[n+1]-x[n]$$
    \textbf{后向差分}为
    $$\nabla x[n]=x[n]-x[n-1]$$
\end{definition}

可以由此定义离散信号的高阶差分，例如，
$$\Delta^2 x[n]=\Delta(\Delta x[n])=x[n+2]-2x[n+1]+x[n]$$

差分方程总能展开为递推式的形式，因此可以仿照微分方程，给出差分方程的定义。其中，我们最关注的是常系数线性差分方程，由这种方程描述的系统是通常是因果的线性时不变系统。
\begin{definition}[常系数线性差分方程]
    关于$y[n]$的差分方程
    $$\sum_{k=0}^{N}a_k y[n-k] = f[n]$$
    称为\textbf{常系数线性差分方程}，其中$a_k$为常数，$f[n]$为已知函数。将$N$称为方程的\textbf{阶数}，将$f[n]$称为方程的\textbf{非齐次项}，如果$f[n]=0$，则称方程为\textbf{齐次差分方程}。
\end{definition}

以上定义中的差分方程是$N$阶的，是因为$N$阶差分能表示为（可以用归纳法证明）：
$$\nabla^N y[n]=\sum_{k=0}^{N}\binom{N}{k}(-1)^{N-k}y[n-k]$$
于是，从$0$到$N$阶差分的线性组合就可以表示为$y[n],y[n-1],\cdots,y[n-N]$的线性组合，即方程左侧的形式。

以上方程是由后向差分方程转化而来的差分方程，即后向差分方程，类似地可以定义前向差分方程，但由于前向差分方程作为系统来分析时，可能不具备因果性，因此我们主要关注后向差分方程。

不难看出，这种方程的解构成线性空间，因此可以将解划分为\textbf{齐次解}与\textbf{特解}两部分。我们仍用记号$y_h[n]$和$y_p[n]$分别表示齐次解与特解。类似于微分方程，齐次解可以由特征方程确定：
\begin{definition}[常系数线性差分方程的特征方程]
    对于常系数线性差分方程
    $$\sum_{k=0}^{N}a_k y[n-k] = f[n]$$
    设解的形式为$y[n]=\lambda^n$，代入可得\textbf{特征方程}为
    $$\sum_{k=0}^{N}a_k \lambda^k = 0$$
\end{definition}

由代数基本定理，特征方程有$N$个复根（重根按重数计算），每个复根对应一个形如$\lambda^n$的齐次解。如果$\lambda$是特征方程的$k$重根，则
$$\lambda^n,n\lambda^n,\dots,n^{k-1}\lambda^n$$
都是齐次解，为了说明这一论断的正确性，我们来验证$n^{k-1}\lambda^n$是齐次解：\begin{align*}
    \sum_{k=0}^{N}a_k (n-k)^{k-1}\lambda^{n-k}&=\sum_{k=0}^{N}a_k \sum_{j=0}^{k-1}\binom{k-1}{j}n^j(-k)^{k-1-j}\lambda^{n-k}\\
    &=\sum_{j=0}^{k-1}\binom{k-1}{j}n^j\lambda^n\sum_{k=0}^{N}a_k (-k)^{k-1-j}\lambda^{-k}\\
    &=\sum_{j=0}^{k-1}\binom{k-1}{j}n^j\lambda^n\frac{1}{j!}\frac{d^{j}}{d\lambda^{j}}\left(\sum_{k=0}^{N}a_k \lambda^{-k}\right)\\
    &=0
\end{align*}
其中，$\sum_{k=0}^{N}a_k \lambda^{-k}$在$\lambda$处的$j$阶导数为零，是因为$\lambda$是特征方程的$k$重根，且$j<k$。

齐次解的线性组合仍为齐次解。一般地，可以将差分方程的齐次解表示为
\begin{align*}
    y_h[n]=\sum_{i=1}^{s}P_i(n) \lambda_i^n=\sum_{i=1}^{s}\sum_{j=0}^{k_i-1}A_{ij} n^j \lambda_i^n
\end{align*}
其中$s$为特征根的个数，$P_i(n)$为关于$n$的$k_i$阶多项式，$k_i$为第$i$个特征根$\lambda_i$的重数，$A_{ij}$为常数。

特解需要根据激励的形式来确定，以下列举几个常见的激励形式及其对应的特解：
\begin{itemize}
    \item 如果$f[n]$为常数，则在$0$不为特征根时，特解为常数；在$0$为$k$重特征根时，特解为关于$n$的$k$次多项式；
    \item 如果$f[n]$为关于$n$的多项式，则特解为关于$n$的同阶多项式；
    \item 如果$f[n]=Ca^n$，其中$a\in\mathbb{R}\setminus \{0\}$不是特征方程的根，则特解为
    $$y_p[n]=Aa^n$$
    其中$A$为常数；
    \item 如果$f[n]=C\sin(\Omega n)$或$f[n]=C\cos(\Omega n)$，其中$\mathrm{e}^{\mathrm{i}\Omega}$不是特征方程的根，则特解为
    $$y_p[n]=A\sin(\Omega n)+B\cos(\Omega n)$$
    其中$A,B$为常数；
    \item 如果$f[n]=P(n)a^n$，其中$a\in\mathbb{R}$是特征方程的$k$重根，且$P(n)$为关于$n$的多项式，则特解为
    $$y_p[n]=n^k Q(n)a^n$$
    其中$Q(n)$为关于$n$的多项式，且$\deg Q=\deg P$；
    \item 如果$f[n]=P(n)\sin(\Omega n)$或$f[n]=P(n)\cos(\Omega n)$，其中$\mathrm{e}^{\mathrm{i}\Omega}$是特征方程的$k$重根，且$P(n)$为关于$n$的多项式，则特解为
    $$y_p[n]=n^k(Q(n)\sin(\Omega n)+R(n)\cos(\Omega n))$$
    其中$Q(n),R(n)$为关于$n$的多项式，且$\max\{\deg Q,\deg R\}=\deg P$。
\end{itemize}
下面给出一个常系数线性差分方程的求解例子：
\begin{example}
    求解差分方程
    $$y[n]-3y[n-1]+2y[n-2]=2^n$$

    解：其特征方程为$\lambda^2-3\lambda+2=0$，有两个不同的根$\lambda_1=1,\lambda_2=2$，因此齐次解为
    $$y_h[n]=A+B\cdot 2^n$$
    设特解的为$y_p[n]=C n 2^n$，代入原方程可得$C=4$，因此特解为$y_p[n]=4 n 2^n$，通解为
    \[y[n]=y_h[n]+y_p[n]=A+B\cdot 2^n+4 n 2^n\]
\end{example}

如果希望消去以上通解中的常数$A$和$B$，就需要两个初始条件。更一般地，一个$N$阶的差分方程的解，需要通过$N$个初始条件来确定。这些初始条件可以通过$n=0$时刻的各阶后向差分来给出。

对于$k$阶后向差分$\nabla^k y[0]$，可以将它展开为$y[-k],y[-k+1],\dots,y[0]$的线性组合，因此给定$n=0$时刻的前$N-1$阶后向差分作为初始条件，就相当于给出
$$y[1-N],y[2-N],\cdots,y[0]$$
这$N$个值作为初始条件。习惯上，我们称这些条件为\textbf{初始值}以区分于微分方程的初始条件。

\subsection{基本解}\label{sub:h_solution}

如果将方程的非齐次项视作系统的激励$e[n]$，将方程的解视作系统的响应$r[n]$，则齐次解对应系统的的零输入响应$y_{zi}[n]$，特解对应系统的零状态响应$y_{zs}[n]$。容易验证，零输入响应对激励具有线性性和时不变性，因此这种系统是一个线性时不变系统，并且我们很快会看到它还是因果的。

对于离散系统，通常认为激励是在$n=0$时刻施加的，此时这个激励就可能影响到上述初始值。连续系统中，我们用起始状态来描述系统的初始储能，从而避开了初值条件依赖于激励的问题。类似地，对于离散系统，定义：
\begin{definition}[离散系统的初始值与零输入初始值]
    设离散系统的激励为$e[n]$，响应为$r[n]$，则系统的\textbf{初始值}为
    $$r[1-N],r[2-N],\cdots,r[0]$$
    系统的\textbf{零输入初始值}为
    $$r_{zi}[1-N],r_{zi}[2-N],\cdots,r_{zi}[0]$$
\end{definition}

初始值可以分解为零输入初始值与零状态初始值（以$n=0$时刻为例）：
\begin{align*}
    r[0]=r_{zi}[0]+r_{zs}[0]
\end{align*}

上一小节中，我们介绍了差分方程的经典解法，为了确定带初始值的差分方程的解，可以先求出含参数的通解，再代入初始值求解这些参数。

但是，将差分方程作为线性时不变系统来分析时，将初始值划分为零输入、零状态初始值更为方便，因为零输入初始值与系统的激励信号无关，是系统的初始储能，代表了系统的历史记忆。

通过离散线性时不变系统的理论，可以将差分方程化为卷积和系统，并通过零输入初始值确定系统的响应，这样就能绕过特解的配凑，而直接得到方程的解。现在，问题又一次回到了对单位脉冲响应$h[n]$的求解，它满足以下方程：
$$\sum_{k=0}^{N}a_k h[n-k] = \delta[n]$$

类比微分方程的情况，设特征方程的根为$\lambda_1,\lambda_2,\cdots,\lambda_s$，重数分别为$k_1,k_2,\cdots,k_s$，则可以猜测单位脉冲响应具有如下形式：
$$h[n]=y_h[n]u[n]=\left(\sum_{i=1}^s P_i(n) \lambda_i^n \right)u[n]$$
其中$P_i(n)$为关于$n$的多项式，$\deg P_i=k_i-1$。

代入$n=0,1,\cdots,N-1$，可以得到$N$个方程，解出多项式$P_i(n)$的系数，即可得到单位脉冲响应$h[n]$，从而得到差分方程的解。

\begin{example}
    求差分方程
    $$y[n]-2y[n-1]+y[n-2]=x[n]$$
    的单位阶跃响应。

    解：特征方程为$\lambda^2-2\lambda+1=0$，有一个重根$\lambda=1$，因此单位脉冲响应的形式为
    $$h[n]=(A n+B)u[n]$$
    代入$n=0,1$，得\begin{align*}
    h[0]=B=1,h[1]=A+B=0
    \end{align*}
    因此$A=-1,B=1$，单位脉冲响应为
    $$h[n]=(1-n)u[n]$$
\end{example}

进一步，还可以用类似的方法讨论以下更一般的差分方程所描述的系统：
\begin{align*}
    \sum_{k=0}^{N}a_k r[n-k] = \sum_{l=0}^{M}b_l e[n-l]
\end{align*}
需要说明的是，当$N\leq M$时，需要在求解单位脉冲响应时考虑$\delta[n]$的后向差分，或更直接地，将单位脉冲响应设为
$$h[n]=\left(\sum_{i=1}^s P_i(n) \lambda_i^n \right)u[n]+\sum_{i=0}^{M-N}B_i\delta[n-i]$$
其中$B_i$为常数，代入$n=0,1,\cdots,M$，解出多项式$P_i(n)$的系数与常数$B_i$，即可得到单位脉冲响应$h[n]$，从而得到差分方程的解。

从基本解的形式可以看出$h[n]$总是因果信号，因此\textbf{差分方程所描述的系统是因果的线性时不变系统}。

在得到单位脉冲响应后，差分方程的解就可以通过卷积和来得到：
\begin{align*}
    r_{zs}[n]=(h*e)[n]=\sum_{m=-\infty}^{\infty}h[n-m]e[m]
\end{align*}
我们又一次将非齐次方程的求解化为卷积的计算。因为这个原因，可以将$h[n]$称为差分方程的\textbf{基本解}。

我们认为激励$e[n]$是在$n=0$时刻施加于系统的，因此$e[n]$和$h[n]$都是因果信号，上式简化为：
\begin{align*}
    r_{zs}[n]=\sum_{m=0}^{n}h[n-m]e[m]
\end{align*}

假设给定了零输入初始值：
\begin{align*}
    r_{zi}[1-N],r_{zi}[2-N],\cdots,r_{zi}[0]
\end{align*}
则零输入响应$r_{zi}[n]$满足齐次方程：
\begin{align*}
    \sum_{k=0}^{N}a_k r_{zi}[n-k] = 0
\end{align*}
上一小节中已经得到齐次方程的通解：
\begin{align*}
    r_{zi}[n]=\sum_{i=1}^{s}P_i(n) \lambda_i^n=\sum_{i=1}^{s}\sum_{j=0}^{k_i-1}A_{ij} n^j \lambda_i^n
\end{align*}
利用零输入初始值，可以得到关于零输入响应$r_{zi}[n]$的$N$个方程，从而解出$r_{zi}[n]$，最终与零状态响应$r_{zs}[n]$相加得到差分方程的解。

综上，通过线性时不变系统的视角，我们得到了与微分方程类似的结果：
\begin{theorem}
    设$h[n]$为差分方程
    $$\sum_{k=0}^{N}a_k y[n-k] = \sum_{l=0}^{M}b_l x[n-l]$$
    的单位脉冲响应，并且有零状态初始值$r_{zs}[1-N],r_{zs}[2-N],\cdots,r_{zs}[0]$，则该差分方程的解为
    \begin{align*}
        y[n]=r_{zi}[n]+r_{zs}[n]=r_{zi}[n]+\sum_{m=0}^{n}h[n-m]x[m]
    \end{align*}
    其中零输入响应$r_{zi}[n]$由零输入初始值唯一确定。
\end{theorem}

沿着\ref{sub:fundamental_solution}的思路，我们尝试由单位阶跃响应来求解单位脉冲响应，从而简化计算。

\begin{lemma}
    设$h[n]$为差分方程
    $$\sum_{k=0}^{N}a_k y[n-k] = \sum_{l=0}^{M}b_l x[n-l]$$
    的单位脉冲响应，$g[n]$为该方程的单位阶跃响应，则$h[n]$是$g[n]$的后向差分：
    $$h[n]=\nabla g[n]=g[n]-g[n-1]$$
\end{lemma}
\begin{proof}
    单位脉冲响应$h[n]$与单位阶跃响应$g[n]$分别满足以下方程：
    \begin{align*}
        \sum_{k=0}^{N}a_k h[n-k] &= \delta[n]\\
        \sum_{k=0}^{N}a_k g[n-k] &= u[n]
    \end{align*}
    将$g[n]$的方程两边同时后向差分，得到\begin{align*}
    \sum_{k=0}^{N}a_k \nabla g[n-k] &= \nabla u[n]=\delta[n]
    \end{align*}
    因此$\nabla g[n]$与$h[n]$满足同一个方程，并且它们的零状态初始值都为零，因此$\nabla g[n]=h[n]$。
\end{proof}

这样，可以少引入一阶$\delta[n]$的差分，从而使方程降一阶。具体来说，对于差分方程
$$\sum_{k=0}^{N}a_k y[n-k] = \sum_{l=0}^{M}b_l x[n-l]$$
当$N\geq M$时，可以将单位阶跃响应设为
$$g[n]=\left(\sum_{i=1}^s P_i(n) \lambda_i^n \right)u[n]$$
其中$P_i(n)$为关于$n$的多项式，$\deg P_i=k_i-1$；当$N<M$时，可以将单位阶跃响应设为
$$g[n]=\left(\sum_{i=1}^s P_i(n) \lambda_i^n \right)u[n]+\sum_{i=0}^{M-N-1}B_i\delta[n-i]$$
其中$B_i$为常数。

\begin{example}
    求差分方程
    $$y[n]-2y[n-1]+y[n-2]=x[n]+x[n-2]$$
    的单位阶跃响应。

    解：特征方程为$\lambda^2-2\lambda+1=0$，有一个重根$\lambda=1$，因此齐次解为
    $$y_h[n]=A n+B$$

    首先，我们直接求解单位脉冲响应。由于$N=M=2$，因此单位脉冲响应的形式为
    $$h[n]=(A n+B)u[n]+C\delta[n]$$
    代入$n=0,1,2$，得\begin{align*}
    \begin{cases}
        h[0]=B+C=1\\
        h[1]=A+B=0\\
        h[2]=2A+B=1
    \end{cases}
    \end{align*}
    因此$A=1,B=-1,C=2$，单位脉冲响应为
    $$h[n]=(n-1)u[n]+2\delta[n]$$

    其次，我们通过求单位阶跃响应来求单位脉冲响应。由于$N=M=2$，因此单位阶跃响应的形式为
    $$g[n]=(A n+B)u[n]$$
    代入$n=0,1$，得\begin{align*}
        \begin{cases}
            h[0]=B=1\\
            h[1]=A+B=0
        \end{cases}
    \end{align*}
    因此$A=-1,B=1$，单位阶跃响应为
    $$g[n]=(1-n)u[n]$$
    从而单位脉冲响应为
    $$h[n]=\nabla g[n]=g[n]-g[n-1]=(n-1)u[n]+2\delta[n]$$
    与之前的结果一致。显然，当单位脉冲响应中含有$\delta[n]$时，\textbf{通过单位阶跃响应求解单位脉冲响应的计算更为简便}。
\end{example}

\subsection{单边z变换}\label{sub:single_z}

根据z变换的移位性质，利用z变换可以将常系数线性差分方程转化为代数方程。为了一并处理差分方程的初始值，我们引入单边z变换：
\begin{definition}[单边z变换]
    离散信号$x[n]$的单边z变换定义为
    \[X(z)=\mathcal{Z}\{x[n]u[n]\}=\sum_{n=0}^{\infty}x[n]z^{-n}\]
\end{definition}

单边z变换相当于因果信号的双边z变换，收敛域为
$$\left\{z\in\mathbb{C}\left||z|>r, r=\varlimsup_{n\to\infty}|x[n]|^{1/n}\right.\right\}$$

单边z变换与双边z变换类似，可以参照\ref{sub:z_operation}给出单边z变换的线性性、移位性质、卷积定理等性质。本小节中，我们简要介绍与双边z变换不同的性质：
\begin{proposition}[单边z变换的移位性质]
    设$x[n]$的单边z变换为$X(z)$，则\begin{align*}
    \mathcal{Z}\{x[n-k]u[n]\}&=z^{-k}X(z)-\sum_{i=0}^{k-1}x[i-k]z^{-i}\\
    \mathcal{Z}\{x[n+k]u[n]\}&=z^{k}X(z)+\sum_{i=-k}^{-1}x[i+k]z^{-i}
    \end{align*}
\end{proposition}
\begin{proof}
    \begin{align*}
        X(z)&=\sum_{i=0}^{\infty}x[i]z^{-i}\\
        &=\sum_{i=k}^{\infty}x[i-k]z^{-(i-k)}\\
        &=z^{k}\sum_{i=k}^{\infty}x[i-k]z^{-i}\\
        &=z^{k}\left[\sum_{i=0}^{\infty}x[i-k]z^{-i}-\sum_{i=0}^{k-1}x[i-k]z^{-i}\right]\\
        &=z^{-k}\left[\mathcal{Z}\{x[n-k]u[n]\}+\sum_{i=0}^{k-1}x[i-k]z^{-i}\right]
    \end{align*}
    即\begin{align*}
    \mathcal{Z}\{x[n-k]u[n]\}&=z^{-k}X(z)-\sum_{i=0}^{k-1}x[i-k]z^{-i}
    \end{align*}
    类似地，\begin{align*}
        X(z)&=\sum_{i=0}^{\infty}x[i]z^{-i}\\
        &=\sum_{i=-k}^{\infty}x[i+k]z^{-(i+k)}\\
        &=z^{-k}\sum_{i=-k}^{\infty}x[i+k]z^{-i}\\
        &=z^{-k}\left[\sum_{i=0}^{\infty}x[i+k]z^{-i}+\sum_{i=-k}^{-1}x[i+k]z^{-i}\right]\\
        &=z^{-k}\left[\mathcal{Z}\{x[n+k]u[n]\}-\sum_{i=-k}^{-1}x[i+k]z^{-i}\right]
    \end{align*}
    即\begin{align*}
    \mathcal{Z}\{x[n+k]u[n]\}&=z^{k}X(z)+\sum_{i=-k}^{-1}x[i+k]z^{-i}
    \end{align*}
\end{proof}

\begin{theorem}[单边z变换的初值定理]
    设$x[n]$的单边z变换为$X(z)$，则\begin{align*}
    x[0]=\lim_{z\to\infty}X(z)
    \end{align*}
\end{theorem}
\begin{proof}
    单边z变换在收敛域内是绝对收敛、内闭一致收敛的，因此可以交换极限与求和顺序，得到\begin{align*}
    \lim_{z\to\infty}X(z)=\lim_{z\to\infty}\sum_{n=0}^{\infty}x[n]z^{-n}=x[0]
    \end{align*}
\end{proof}

\begin{theorem}[单边z变换的终值定理]
    设$x[n]$的单边z变换为$X(z)$，如果$X(z)$在$z=1$处有简单极点，并且$X(z)$的其余奇点的模都小于1，则\begin{align*}
    \lim_{n\to\infty}x[n]=\lim_{z\to 1}(z-1)X(z)
    \end{align*}
\end{theorem}
\begin{proof}
    $X(z)$的收敛域为$|z|>1$。由逆z变换的留数定理公式，以半径$R>1$的圆$C_R:t\mapsto Re^{\mathrm{i}t},t\in[0,2\pi]$为围道，有\begin{align*}
    x[n]u[n]=\frac{1}{2\pi\mathrm{i}}\oint_{C_R}X(z)z^{n-1}\mathrm{d}z=\mathrm{Res}(X(z)z^{n-1},1)+\sum_{k}\mathrm{Res}(X(z)z^{n-1},p_k)
    \end{align*}
    其中$p_k$为$X(z)$的其余极点，$|p_k|<1$。
    \begin{align*}
        \mathrm{Res}(X(z)z^{n-1},1)&=\lim_{z\to 1}(z-1)X(z)z^{n-1}=\lim_{z\to 1}(z-1)X(z)
    \end{align*}
    为了处理它们的留数，可以选取一个半径$r$，使得$|p_k|<r<1$，以半径$r$的圆$C_{r}:t\mapsto re^{\mathrm{i}t},t\in[0,2\pi]$为围道，有\begin{align*}
    \left|\sum_{k}\mathrm{Res}(X(z)z^{n-1},p_k)\right|&=\left|\frac{1}{2\pi\mathrm{i}}\oint_{C_r}X(z)z^{n-1}\mathrm{d}z\right|\\
    &\leq \frac{1}{2\pi}\oint_{C_r}|X(z)|r^{n-1}\mathrm{d}t\\
    &\leq \max_{|z|=r}|X(z)|r^{n-1}
    \end{align*}
    令$n\to\infty$，由于$r<1$，上式以指数函数的速度趋于$0$，因此\begin{align*}
    \lim_{n\to\infty}x[n]&=\frac{1}{2\pi\mathrm{i}}\lim_{n\to\infty}\int_{C_R}X(z)z^{n-1}\mathrm{d}z\\
    &=\lim_{n\to\infty}\mathrm{Res}(X(z)z^{n-1},1)+\lim_{n\to\infty}\sum_{k}\mathrm{Res}(X(z)z^{n-1},p_k)\\
    &=\lim_{z\to 1}(z-1)X(z)
    \end{align*}
\end{proof}

我们对离散信号的终值定理做三点说明。

首先，通过将$X(z)$限制在单位圆上，可以由这个定理得到离散时间傅里叶变换的终值定理，因此可以用“\textbf{频域上的低频成分对应时域的稳态行为}”来解释这个定理。我们知道，对于连续系统，终值定理为$\lim_{t\to\infty}f(t)=\lim_{s\to 0}sF(s)$，与离散系统一样，它们的终值定理都是向变换域上的$0$频率点取极限。

其次，终值定理的这种证明方法只适用于$H(z)$无本性奇点的情况，在\ref{sub:phase_response}中我们曾使用过类似的方法。当然，也可以仿照\ref{sub:initial_and_final_value}中的思路，将$X(z)$分解为$z=1$处的一阶分式与$|z|>r$上的全纯函数的和，这个方法在$H(z)$的奇点不全为极点时同样适用。

最后，在多数文献中，这个定理同样允许$z=1$不是奇点，但根据洛朗级数的收敛半径公式，当奇点全部位于单位圆内（不含边界）时，可以判断$X(z)$的收敛半径小于$1$，\textbf{从而$|x[n]|$的上极限为$0$}，终值定理仍然成立，但是我们仅需考虑非平凡情况。

进一步，当单位圆外有奇点时，还可以判断$|x[n]|$的上极限为$\infty$，\textbf{$x[n]$无界但极限未必存在}，此时终值定理不成立。可见，类似于\ref{sub:converge_abscissa}，我们还是\textbf{能够根据$X(z)$在$z$平面上的奇点情况来判断原信号的时域特性}。

\subsection{常系数线性差分方程的单边z变换解法}\label{sub:z_difference_equation}

考虑以下差分方程：
$$\sum_{k=0}^{N}a_k y[n-k] = \sum_{l=0}^{M}b_l x[n-l]$$


两边同时乘$u[n]$并做单边z变换：
\begin{align*}
    \sum_{k=0}^{N}a_k \mathcal{Z}\{y[n-k]u[n]\}&=\sum_{l=0}^{M}b_l \mathcal{Z}\{x[n-l]u[n]\}
\end{align*}
利用上一小节中得到的移位性质，有\begin{align*}
    \sum_{k=0}^{N}a_k \left(z^{-k}Y(z)-\sum_{i=0}^{k-1}y[i-k]z^{-i}\right)&=\sum_{l=0}^{M}b_l \left(z^{-l}X(z)-\sum_{j=0}^{l-1}x[j-l]z^{-j}\right)
\end{align*}
其中，\begin{align*}
    k=0&\implies \sum_{i=0}^{k-1}y[i-k]z^{-i}=0\\
    l=0&\implies \sum_{j=0}^{l-1}x[j-l]z^{-j}=0
\end{align*}
因此，\begin{align*}
    Y(z)=\frac{\sum_{l=0}^{M}b_l z^{-l}X(z)+\sum_{k=1}^{N}\sum_{i=0}^{k-1}a_k y[i-k]z^{-i}-\sum_{l=1}^{M}\sum_{j=0}^{l-1}b_l x[j-l]z^{-j}}{\sum_{k=0}^{N}a_k z^{-k}}
\end{align*}

特别地，当方程右边恰为$x[n]$时，$M=0,b_0=1$，上式简化为\begin{align*}
    Y(z)=\frac{X(z)+\sum_{k=1}^{N}\sum_{i=0}^{k-1}a_k y[i-k]z^{-i}}{\sum_{k=0}^{N}a_k z^{-k}}
\end{align*}

\begin{example}
    求解差分方程
    $$y[n]-3y[n-1]+2y[n-2]=2^n$$
    的解，初始值为$y[-2]=1,y[-1]=0$。

    解：将方程两边同时乘$u[n]$并做单边z变换，得到\begin{align*}
    Y(z)-3\left(z^{-1}Y(z)-y[-1]\right)+2\left(z^{-2}Y(z)-y[-1]z^{-1}-y[-2]\right)=\frac{z}{z-2}
    \end{align*}
    代入并整理得\begin{align*}
    \left(1-3z^{-1}+2z^{-2}\right)Y(z)&=\frac{z}{z-2}-2
    \end{align*}
    因此\begin{align*}
    Y(z)=\frac{4-z}{z-2}\cdot\frac{1}{1-3z^{-1}+2z^{-2}}=\frac{z^2(4-z)}{(z-2)^2(z-1)}
    \end{align*}
    由逆z变换的留数定理公式，选取$|z|>2$中的围道$C$，有\begin{align*}
    y[n]u[n]&=\frac{1}{2\pi\mathrm{i}}\oint_{C}Y(z)z^{n-1}\mathrm{d}z\\
    &=\mathrm{Res}\left(\frac{z^2(4-z)}{(z-2)^2(z-1)}z^{n-1},2\right)+\mathrm{Res}\left(\frac{z^2(4-z)}{(z-2)^2(z-1)}z^{n-1},1\right)\\
    &=\lim_{z\to 2}\frac{d}{dz}\left(\frac{z^{n+1}(4-z)}{(z-1)}\right)+\lim_{z\to 1}\frac{z^2(4-z)}{(z-2)^2}z^{n-1}\\
    &=\lim_{z\to 2}\frac{[(n+1)z^n(4-z)-z^{n+1}](z-1)-z^{n+1}(4-z)}{(z-1)^2}+\lim_{z\to 1}\frac{z^{n+1}(4-z)}{(z-2)^2}\\
    &=(2n-4)2^n+3,n\geq 0
\end{align*}
\end{example}
